\subsection{Odporúčací systém}                                                                                                                                                                                                   
                                                                                          
  Pre aplikáciu Joinly sme zvolili Weighted Hybrid System -
  váhový hybridný systém kombinujúci tri prístupy:
                                          
  \begin{itemize}[noitemsep, nolistsep]                                                                                                                                                                                                                 
      \item \textbf{Content-Based} — odporúčania na základe kategórií                                                                                                                                                                          navštívených eventov,                                                                                                                                                                                                  
      \item \textbf{Kontextový} — zohľadňuje vzdialenosť, čas a popularitu,                                                                                                                                                        
      \item \textbf{Explicitná spätná väzba} — bonus za eventy z obľúbených
            kategórií.
  \end{itemize}

  Každý z týchto prístupov prispieva do výsledného skóre s nastaviteľnou
  váhou, ktorú si používateľ môže upraviť priamo v aplikácii pomocou sliderov.

  \paragraph{Prečo hybridný systém?}
  \begin{enumerate}[noitemsep, nolistsep]
      \item \textbf{Kvalitné metadáta} — každý event má kategóriu,
            podkategóriu, lokáciu a hodnotenie.
      \item \textbf{Menšia používateľská báza} — Collaborative filtering
            vyžaduje veľké množstvo používateľov; hybridný prístup funguje
            aj bez nich.
      \item \textbf{Vysvetliteľnosť} — systém dokáže zdôvodniť odporúčanie
            (napr.\ „odporúčame ti toto, lebo máš rád športové eventy").
      \item \textbf{Flexibilita} — používateľ si sám nastaví čo je pre neho
            dôležitejšie (blízkosť vs.\ kategória vs.\ popularita).
      \item \textbf{Real-time výpočet} — skóre sa počíta za behu,
            nevyžadujú sa predpočítané modely.
  \end{enumerate}

  \paragraph{Dvojúrovňové kategórie:}
  Štandardný Content-Based systém používa jednu úroveň kategórií.
  V našom systéme sú kategórie dvojúrovňové:
  \begin{enumerate}[noitemsep, nolistsep]
      \item \textbf{Hlavná kategória} — široká skupina (napr.\ \textit{Sports}),
      \item \textbf{Podkategória} — špecifická (napr.\ \textit{Cycling sports}).
  \end{enumerate}

  Toto umožňuje presnejšie odporúčania — používateľ, ktorý navštívil
  cyklistický event, dostane bonus pre ďalšie cyklistické eventy, nie len
  všeobecne športové. Previazanie medzi úrovňami zabezpečuje trieda
  \textbf{CategoryMapper}.

  % ═══════════════════════════════════════════════════════════════════════                                                                                                                                                        
  \subsubsection{Váhy a výsledné skóre}
  % ═══════════════════════════════════════════════════════════════════════                                                                                                                                                        

  Systém používa 6 faktorov, každý s vlastnou váhou. Predvolené hodnoty
  sú definované v dátovej triede \texttt{RecommendationWeights}:

  \begin{lstlisting}[language=Kotlin, caption=Predvolené váhy odporúčacieho systému]
  data class RecommendationWeights(
      val mainCategory  : Float = 0.20f,  // 20 %                                                                                                                                                                                  
      val subCategory   : Float = 0.16f,  // 16 %                                                                                                                                                                                  
      val distance      : Float = 0.16f,  // 16 %                                                                                                                                                                                  
      val rating        : Float = 0.16f,  // 16 %                                                                                                                                                                                  
      val popularity    : Float = 0.16f,  // 16 %                                                                                                                                                                                  
      val favoriteBonus : Float = 0.16f   // 16 %                                                                                                                                                                                  
  )
  \end{lstlisting}

  Hlavná kategória má mierne vyššiu váhu (20\,\%) pretože je
  najspoľahlivejším indikátorom záujmu používateľa.
   Keďže súčet váh musí byť vždy rovný 1 (100\,\%), systém hodnoty
  automaticky normalizuje po každej zmene slidera.

  \begin{equation}
  \sum_{i=1}^{6} w_i = 1{,}0 \quad \text{(vždy 100\,\%)}
  \end{equation}

  \paragraph{Vzorec celkového skóre:}

  \begin{equation}
  \boxed{
  Score = w_1 \cdot C_{main} + w_2 \cdot C_{sub} + w_3 \cdot D +
          w_4 \cdot R + w_5 \cdot P + w_6 \cdot B_{fav}
  }
  \end{equation}

  Kde:
  \begin{itemize}[noitemsep, nolistsep]
      \item $w_1 \ldots w_6$ — nastaviteľné váhy (súčet = 1,0),
      \item $C_{main}$ — skóre hlavnej kategórie $(0{,}0 - 1{,}0)$,
      \item $C_{sub}$ — skóre podkategórie $(0{,}0 - 1{,}0)$,
      \item $D$ — skóre vzdialenosti $(0{,}0 - 1{,}0)$,
      \item $R$ — skóre ratingu $(0{,}0 - 1{,}0)$,
      \item $P$ — skóre popularity $(0{,}0 - 1{,}0)$,
      \item $B_{fav}$ — bonus za obľúbenú kategóriu $(0{,}0 - 1{,}0)$.
  \end{itemize}

  Keďže každý faktor je v rozsahu $0{,}0 - 1{,}0$ a súčet váh je vždy
  1,0, výsledné skóre je garantovane v rozsahu:

  \begin{equation}
  0{,}0 \leq Score \leq 1{,}0
  \end{equation}

  % ═══════════════════════════════════════════════════════════════════════                                                                                                                                                        
  \subsubsection{Jednotlivé faktory}
  % ═══════════════════════════════════════════════════════════════════════                                                                                                                                                        

  \paragraph{Hlavná kategória ($C_{main}$):}
  Systém zostaví profil používateľa z navštívených eventov - koľko percent
  navštívených eventov patrilo do ktorej kategórie. Skóre sa vypočíta
  nasledovne:

  \begin{lstlisting}[language=Kotlin, caption=Výpočet skóre hlavnej kategórie]
  var mainMatch = userProfile[mainCategory] ?: 0.0
  if (userProfile.isEmpty()) mainMatch = 0.5   // novy pouzivatel
  else if (mainMatch == 0.0) mainMatch = 0.05  // nikdy nenavstivena

  val mainScore = mainMatch * w.mainCategory
  \end{lstlisting}

  \begin{itemize}[noitemsep, nolistsep]
      \item Ak používateľ navštívil 10 eventov, z toho 4 boli „Sports"
            $\rightarrow$ preferencia Sports = 40\,\% $(C_{main} = 0{,}4)$.
      \item Pre nových používateľov (bez histórie) $\rightarrow$                                                                                                                                                                   
            $C_{main} = 0{,}5$ (neutrálne skóre).
      \item Pre nikdy nenavštívenú kategóriu $\rightarrow$                                                                                                                                                                         
            $C_{main} = 0{,}05$ (aby systém nevrátil prázdny zoznam).
  \end{itemize}

  \paragraph{Podkategória ($C_{sub}$):}
  Funguje rovnako ako hlavná kategória, ale na úrovni podkategórií.
  Skóre podkategórie dopĺňa hlavnú kategóriu a umožňuje rozlišovať
  medzi napríklad cyklistickými a plaveckými eventmi v rámci rovnakej
  skupiny Sports.

  \paragraph{Vzdialenosť ($D$):}
  Vzdialenosť sa počíta Haversinovým vzorcom (popísaným v sekcii
  filtrovania) a normalizuje sa tak, aby bližšie eventy dostávali
  vyššie skóre:

  \begin{lstlisting}[language=Kotlin, caption=Výpočet skóre vzdialenosti]
  val distance  = haversineKm(userLat, userLng,
                               event.latitude, event.longitude)
  val distScore = (1.0 / (1.0 + distance / 10.0)) * w.distance
  \end{lstlisting}

  \paragraph{Rating ($R$):}
  Hodnotenie eventu sa normalizuje do rozsahu $0{,}0 - 1{,}0$. Eventy
  bez hodnotení dostanú neutrálne skóre 0,5, čím sa zabraňuje ich
  nespravodlivému penalizovaniu:

  \paragraph{Popularita ($P$):}
  Popularita sa meria ako pomer aktuálneho počtu účastníkov voči
  maximálnej kapacite eventu. Ak event nemá stanovený limit, používa
  sa strop 100:


  \paragraph{Bonus za obľúbenú kategóriu ($B_{fav}$):}
  Obľúbené eventy sa \textbf{nepočítajú} do profilu preferencií —
  používajú sa \textbf{výhradne} na tento bonus. To umožňuje
  používateľovi explicitne označiť čo sa mu páči bez ovplyvnenia
  celkového profilu navštívených eventov:

  % ═══════════════════════════════════════════════════════════════════════                                                                                                                                                        
  \subsubsection{Výstup a používateľské rozhranie}
  % ═══════════════════════════════════════════════════════════════════════                                                                                                                                                        

  Výsledkom výpočtu je zoznam objektov ScoredEvent zoradený
  zostupne podľa skóre. Odporúčania sa zobrazujú na obrazovke
  RecommendedEventsScreen, kde každý event zobrazuje aj
  percentuálne skóre zhody. Používateľ si môže jednotlivé váhy faktorov
  upraviť pomocou sliderov priamo na tejto obrazovke — po každej zmene
  sa odporúčania prepočítajú v reálnom čase.